\chapter{Introducción}

\section{Resumen}

En el presente documento se presenta la primera etapa de una investigación orientada a la simulación de sistemas complejos en espacios bidimensionales mediante autómatas celulares neuronales. El objetivo principal es utilizar una red neuronal para definir la función de transición de un autómata celular, con el fin de modelar y simular las representaciones de sistemas complejos sobre una red regular bidimensional.

Los resultados incluyen el desarrollo del marco teórico, que recopila el conjunto de definiciones formales necesarias; el estado del arte, donde se analizan investigaciones previas relacionadas con diferentes enfoques y configuraciones; la descripción del marco metodológico, en donde se especifican las técnicas, herramientas y procedimientos utilizados; y los avances obtenidos hasta el momento, que abarcan la implementación y documentación de herramientas para la simulación de autómatas celulares en dos dimensiones.

\section{Planteamiento del problema}

La simulación de sistemas complejos ha surgido como una herramienta indispensable para la comprensión y el análisis de fenómenos en diversas áreas del conocimiento, desde la física y la biología hasta las ciencias sociales y económicas. Tradicionalmente, los autómatas celulares han servido como un modelo matemático fundamental para estudiar estos sistemas mediante la representación de procesos dinámicos. 

Los autómatas celulares son un modelo matemático definido como una 4-tupla que contiene: 
\begin{itemize}
	\item Un espacio de celdas en un dominio discreto.
	\item Un conjunto de estados posibles.
	\item Una vecindad que define las dependencias entre células.
	\item Una función de transición que determina cómo evolucionan simultáneamente las células\cite{wolframCA, ZenilMartinez2024}. 
\end{itemize}

Ejemplos clásicos de autómatas celulares incluyen el Juego de la Vida de John Conway y los 256 autómatas elementales descritos por Stephen Wolfram. Estas herramientas han sido aplicadas en múltiples contextos como la simulación de incendios, modelado de distribuciones poblacionales, simulación de partículas y el análisis de ciclos climáticos \cite{darg}. No obstante, una limitación crítica de los modelos convencionales reside en la rigidez de sus funciones de transición; estas deben definirse de manera estática y anticipada, restringiendo la posibilidad de adaptarse a condiciones imprevistas o cambios en el entorno. Esta limitación motiva la búsqueda de nuevas metodologías capaces de incorporar mayor flexibilidad y adaptabilidad en la simulación de sistemas dinámicos. 

En este contexto, se busca la integración de redes neuronales en el marco de los autómatas celulares. Las redes neuronales, inspiradas en el funcionamiento del cerebro y su capacidad de aprendizaje, ofrecen la posibilidad de redefinir las reglas de transición de forma dinámica. Esto permite que el autómata celular se ajuste de manera autónoma a partir de la experiencia y el entrenamiento \cite{RN}. 

El presente proyecto de investigación propone superar las limitaciones de los autómatas celulares mediante la integración de redes neuronales en la definición de las reglas de transición, dando origen a los \textit{autómatas celulares neuronales}. En este enfoque, las reglas de actualización dejan de ser fijas y adquieren la capacidad de aprender y adaptarse a partir de los datos, lo que permite que el sistema se ajuste de forma dinámica ante comportamientos complejos y entornos variables \cite{redesNeuronales, distill}. Esta aproximación no solo busca ampliar el campo de aplicación de los autómatas celulares, sino que también busca abrir nuevas oportunidades para modelar y simular fenómenos de formas realistas.  

Dadas las infinitas configuraciones de autómatas celulares, este proyecto se limita a un espacio bidimensional finito (con condiciones de frontera) y un conjunto finito y numerable de estados. Adicionalmente, se busca experimentar con diferentes vecindades incluyendo Moore y Newman. Aunque el objetivo no es determinar las condiciones para la evolución de las células, si se busca realizar pruebas bajo diferentes configuraciones que permitan mostrar los resultados sin comprometer su calidad.

Por otro lado, la red neuronal recibe como parámetros los estados de las celdas vecinas y devuelve el nuevo estado de la celda a la cual se aplique la función de transición. Este proceso es el mismo para todas las células.

Es por esto que surge la pregunta central de esta investigación: ¿Cómo los autómatas celulares neuronales pueden simular sistemas complejos en un espacio bidimensional? A partir de este cuestionamiento, surgen las siguientes preguntas de investigación.

\subsection{Preguntas de investigación}

\begin{itemize}
	\item ¿Qué es un sistema complejo?
	\item ¿Qué es simular?
	\item ¿Como se simula un sistema complejo?
	\item ¿Qué es un  autómata celular?
	\item ¿Que es una red neuronal?
	\item ¿Qué es un autómata celular neuronal?
	\item ¿Como simular un sistema complejo?
\end{itemize}

\section{Objetivo}

Desarrollar un sistema de simulación basado en autómatas celulares neuronales bidimensionales que permita analizar la generación, mantenimiento, recuperación y propagación de patrones de información en entornos complejos y dinámicos.

\subsection{Objetivos específicos}

\begin{itemize}
	\item Diseñar e implementar un simulador de espacios discretos bidimensionales que sirva de plataforma para estudiar los autómatas celulares neuronales.
	\item Entrenar y evaluar el desempeño de los autómatas celulares neuronales bajo diferentes configuraciones.
	\item Replicar configuraciones clásicas con comportamientos y soluciones conocidas. 
	\item Analizar y comparar diferentes comportamientos con base en la red neuronal establecida.
\end{itemize}

\section{Hipótesis}

Si un sistema de simulación está basado en autómatas celulares neuronales bidimensionales, entonces es posible analizar y reproducir la generación, mantenimiento y propagación de patrones de información en entornos complejos y dinámicos.

\section{Justificación}

Los autómatas celulares son modelos lo bastante simples para un análisis teórico, pero lo suficientemente complejos para exhibir fenómenos emergentes. Definir las reglas de comportamientos erráticos para ser sincronizados resulta en una tarea compleja \cite{computationalBiology2016}. 

Por este motivo, el presente proyecto de investigación surge ante la necesidad de superar las limitaciones inherentes a los modelos tradicionales de autómatas celulares, los cuales se caracterizan por depender de reglas de transición fijas y predeterminadas. Este enfoque restrictivo, si bien ha sido útil para modelar ciertos fenómenos, no siempre captura la complejidad y la adaptabilidad de sistemas dinámicos presentes en entornos naturales, sociales y tecnológicos. 

Cada autómata celular resuelve un problema, sin embargo, su definición matemática implica que existe una infinitud de problemas a resolver. Por lo que realizar configuraciones hasta obtener el resultado deseado resulta un proceso computacionalmente inviable \cite{elYacoubi2021control}. 

En este sentido, se propone la integración de redes neuronales en la definición de las reglas de transición, con el objetivo de desarrollar un modelo de \textit{autómata celular neuronal} capaz de aprender y adaptarse de manera autónoma a cambios en el entorno. Esto no solo amplía el alcance y la aplicabilidad de los autómatas celulares, sino que también abre un novedoso campo de estudio en el que se combinan herramientas de inteligencia artificial y simulación de sistemas complejos. 
La principal innovación de este enfoque reside en reemplazar las funciones de transición estáticas por modelos ajustables. Mediante técnicas de aprendizaje automático, la función de transición puede entrenarse con datos representativos, permitiendo que el sistema se adapte de manera autónoma a los cambios en el entorno \cite{pre}. 

Adicionalmente, este proyecto profundiza en técnicas de simulación, modelado matemático y redes neuronales, integrando conocimientos de ciencias de la computación, vida artificial y matemáticas.

La experiencia adquirida durante el desarrollo e implementación de estos modelos abrirá puertas a futuras investigaciones y proyectos en el campo de la inteligencia artificial. Además, la creación de un entorno de código abierto contribuirá al avance del estado del arte y facilitará que otros investigadores puedan extender y aplicar estos conceptos.

En síntesis, este proyecto busca contribuir al avance de la simulación de sistemas complejos mediante la creación de un modelo flexible, adaptable y extensible de autómatas celulares neuronales, consolidando un punto de encuentro entre la teoría de sistemas dinámicos y las técnicas modernas de aprendizaje automático.

\section{Alcances y limitaciones}

\textbf{Alcances}

\begin{itemize}
	\item El estudio se limita a la simulación en espacios bidimensionales finitos con condiciones de frontera establecidas en una red regular definida.
	\item Se emplean redes neuronales supervisadas para la definición de la función de transición.
	\item Se analizarán comportamientos de patrones conocidos y emergentes bajo distintas configuraciones de estados y vecindades.
\end{itemize}

\textbf{Limitaciones}

\begin{itemize}
	\item No se contempla la generalización del modelo a diferentes dimensionalidades o topologías de red.
	\item El análisis se centra en la simulación de sistemas complejos, sin abordar la optimización computacional del entrenamiento.
	\item Los resultados se interpretan únicamente dentro del contexto de los parámetros y condiciones experimentales definidos.
\end{itemize}
